\selectlanguage{ngerman}%
\begin{ledgroupsized}[r]{120mm}%
\footnotesize%
\pstart%
\noindent\textbf{Überlieferung:}%
\pend%
\end{ledgroupsized}%
\begin{ledgroupsized}[r]{114mm}%
\footnotesize%
\pstart%
\parindent -6mm%
\makebox[6mm][l]{\textit{L}}%
Aufzeichnung:
LH~XXXVII~4, Bl.\;87.
Ein Blatt~4\textsuperscript{o},
Ränder teilweise beschnitten mit geringfügigem Textverlust;
Papiererhaltungsmaßnahmen.
\href{\Scan{LH_37_04_087wz}}%
Fragment eines {Wasserzeichens}:
Harzer Papier.
Eine Seite auf Bl.\;\href{\Scan{LH_37_04_087r}}{87\,r\textsuperscript{o}}.
Bl.\;\href{\Scan{LH_37_04_087v}}{87\,v\textsuperscript{o}} ist unbeschrieben.
\pend%
\end{ledgroupsized}%
\vspace{5mm}%
\begin{ledgroup}%
\footnotesize%
\pstart%
\noindent%
\textbf{Datierungsgründe:}
Die Datierung der Aufzeichnung N.\;\ref{60243} beruht auf äußeren Eigenschaften des Text\-trägers %
sowie auf inhaltlichen Betrachtungen.
Siehe die ausführliche Begründung auf S.\;\pageref{60243_dat}.
\pend%
\end{ledgroup}%
\selectlanguage{latin}%
\frenchspacing%
\vspace{8mm}%
\normalsize%
\pstart%
\noindent%
\lbrack%
\href{\Scan{LH_37_04_087r}}%
{87\,r\textsuperscript{o}}%
\rbrack\ %
Si corpus insequens%
\indexS{corpus insequens}
impingat in tardius antecedens,%
\indexS{corpus antecedens}
quaeritur
quanta vi in ipsum agere%
\indexS{vis agendi}
intelligatur,
an ea celeritatis suae parte %
qua ipsius celeritatem vincit,
an vero celeritate sua tota?
Verum sciendum est
neutrum proprie locum habere,
nec in aestimanda vi percussionis%
\indexS{vis percussionis aestimanda}
rationem habendam celeritatis absolutae%
\indexS{celeritas absoluta}
sed respectivae,%
\indexS{celeritas respectiva}
qua scilicet corpora sibi accedunt.
\edtext{%
Et cum arbitrarium sit
quantanam alterutri sit tribuenda celeritas,%
}{%
\lemma{Et}%
\Bfootnote{%
\textit{(1)}~cum diversa futura essent phaenomena %
\textit{(2)}~cum arbitrarium %
\lbrack...\rbrack\ tribuenda celeritas,%
~\textit{L}%
}}
seu quomodo celeritates accessionum%
\indexS{celeritas accessionis}%
\indexS{accessio}
inter ipsa corpora sint partiendae,%
\indexS{celeritas partienda}
\lbrack\textendash\rbrack\
potest enim unum fingi quiescens %
aut tarde motum,
toto motu aut maxima ejus parte in
\edtext{%
alterum transcripta,%
\indexS{motus transcriptus in aliud corpus}%
}{%
\lemma{alterum}%
\Bfootnote{%
\textit{(1)}~transcripto
\textit{(2)}~transcripta%
~\textit{L}%
}}
\lbrack\textendash\rbrack\
ideo calculum ita esse instituendum,%
\indexS{calculus de potentia respectiva}
ut celeritates accessionum%
\indexS{celeritas accessionis}
sint
\edtext{%
corporibus reciproce proportionales.%
}{%
\lemma{corporibus}%
\Bfootnote{%
\textit{(1)}~proportionales %
\textit{(2)}~reciproce proportionales.%
~\textit{L}%
}}
Itaque
\edlabel{LH_37_04_087_087r_potabs-1}%
licet potentia absoluta%
\indexS{potentia absoluta aestimanda}
sit aestimanda ductu corporis in quadratum celeritatis,%
\indexS{ductus corporis in quadratum celeritatis}
potentia tamen respectiva%
\indexS{potentia respectiva aestimanda}
aestimabitur ductu corporis in celeritatem,%
\indexS{ductus corporis in celeritatem}
quando de ipsa inter
\edtext{%
corpora concurrentia partienda%
\indexS{potentia respectiva partienda}%
}{%
\lemma{corpora}%
\Bfootnote{%
\textit{(1)}~inter se partienda %
\textit{(2)}~concurrentia partienda%
~\textit{L}%
}}
agitur.%
\edlabel{LH_37_04_087_087r_potabs-2}
Cujus rei ratio ex eo fluit,%
\indexS{ratio fluens}
quod centrum
\edtext{gravitatis compositi%
\indexS{centrum gravitatis compositi}
ex corporibus %
necesse est eodem%
}{%
\lemma{gravitatis}%
\Bfootnote{%
\textit{(1)}~necesse est eo %
\textit{(2)}~compositi ex %
\lbrack...\rbrack\ est eodem%
~\textit{L}%
}}
semper modo progredi.
\pend%
\pstart%
Ponamus duas
\edtext{%
Tabulas%
\indexS{tabula affixa librae}
\textit{LP} et \textit{QM}%
}{%
\lemma{Tabulas \textit{LP} et \textit{QM}}%
\Cfootnote{%
Vgl. das Diagramm \bracks{\textit{Fig.\:1}}.%
}}
horizonti parallelas
esse brachiis \textit{EL} et \textit{EM}%
\indexS{brachium librae}
librae \textit{LEM} %
\indexS{libra}
circa centrum \textit{E}%
\indexS{centrum librae}
mobilis
\edtext{%
affixas,
et in tabula \textit{PL}%
\indexS{tabula affixa librae}
esse duos globulos \textit{A} et \textit{B},%
\indexS{globulus}
in Tabula autem \textit{MQ}%
\indexS{tabula affixa librae}
duos alios \textit{C} et \textit{D},
omn\lbrack ia\rbrack\
ita sita et constituta %
ut tabulae tam%
\indexS{tabula affixa librae}
sine globulis
quam etiam cum globulis
sint in aequilibr\textlangle io\textrangle%
\indexS{aequilibrium}
et%
}{%
\lemma{affixas}%
\Bfootnote{%
\bars{et in aequilibrio constitutas \textit{erg.~u. gestr.}}%
, et in %
\lbrack...\rbrack\ \textit{C} et \textit{D}, %
\bars{omnes \textit{ändert Hrsg.}} ita 
\textit{(1)}~sitos %
\textit{(2)}~sita %
\textit{(a)}~ut duae tabulae sint in aequilibrio et %
\textit{(b)}~et %
\textit{(aa)}~\textlangle ipsi vim\textrangle\ %
\textit{(bb)}~constituta ut %
\lbrack...\rbrack\ aequilibr\textlangle io\textrangle\ et%
~\textit{L}%
}}
\ledpb%
maneant adhuc in aequilibrio%
\indexS{aequilibrium} %
licet globuli%
\indexS{globulus}
\edtext{%
a\lbrack d\rbrack%
}{%
\lemma{}%
\Bfootnote{%
a~%
\textit{L~ändert Hrsg.}%
}}
se invicem accedant %
aut a se invicem recedant, %
aut etiam %
ubi ad extremum Tabulae%
\indexS{tabula affixa librae}
accessere %
inde reflectantur %
eadem
(ut suppono)
celeritate
qua
\edtext{%
venerunt.
His positis si%
}{%
\lemma{venerunt.}%
\Bfootnote{%
\textit{(1)}~Si %
\textit{(2)}~His positis si%
~\textit{L}%
}}
jam intelligantur globi duo \textit{A} et \textit{B}%
\indexS{globi occurrentes sibi directe}
ad se accedentes,
sibi directe occurrere %
atque pergere vel repercuti,
globuli autem \textit{C} et \textit{D}%
\indexS{globulus}
sese praeterire,
nihilo minus tam ante %
quam post concursum%
\indexS{concursus corporum}
debebit manere
aequilibrium,%
\indexS{aequilibrium}
unde necesse est 
centrum gravitatis%
\indexS{centrum gravitatis commune}
globulorum%
\indexS{globulus}
unius partis
eodem modo procedere
ut centrum gravitatis%
\indexS{centrum gravitatis commune}
globulorum%
\indexS{globulus}
alterius partis,
seu duo illa centra
\edtext{%
gravitatis semper a centro librae \textit{E}%
\indexS{centrum librae}
eodem modo recedere,
vel eodem modo ad ipsum accedere,
alioqui nunc uno nunc altero centrorum gravitatis%
\indexS{centrum gravitatis commune}
alternis accedente %
magis minusve ad
centr\lbrack um\rbrack\
librae%
\indexS{centrum librae}%
\indexS{libra}
vel ab eo
\lbrack recedente\rbrack,
nunc unum nunc%
}{%
\lemma{gravitatis}%
\Bfootnote{%
\textit{(1)}~eandem semper distantiam servare %
\textit{(2)}~semper a %
\lbrack...\rbrack\ accedere, alioqui
\textit{(a)}~altern\textlangle is\textrangle\ %
\textit{(b)}~nunc uno nunc altero %
\bars{centrorum gravitatis \textit{erg.}}
alternis accedente %
\textit{(aa)}~ac %
\textit{(bb)}~\bars{vel recedente \textit{streicht Hrsg.}} magis minusve ad %
\bars{centro \textit{ändert Hrsg.}} librae vel ab eo %
\bars{recedente \textit{erg. Hrsg.}} nunc unum nunc%
~\textit{L}%
}}
alterum brachium librae%
\indexS{libra}%
\indexS{brachium librae}
inclinabitu\textlangle r\textrangle\
certa vi,%
\indexS{vis inclinationis}
habebiturque motus perpetuus mechanicus.%
\indexS{motus perpetuus mechanicus}
\pend%
\vspace{3.00em} %
\centerline{%
\includegraphics[width=0.35\textwidth]{%
gesamttex/edit_VIII-4/images/LH_37_04_087_d_087r.pdf%
}}
\nobreak\vspace{0.5em}%
\centerline{%
\lbrack\textit{Fig.\:1}\rbrack%
}
\vspace{5.50em}%
\pstart%
\noindent%
\centering%
*    *    *
\pend%
\newpage%
\selectlanguage{ngerman}%
\begin{ledgroup}%
\footnotesize%
\pstart%
\noindent%
\textls{Ausführliche Datierungsgründe:}%
\pagelabel{60243_dat}
Der Träger der titellosen und undatierten Aufzeichnung N.\;\ref{60243} weist das Fragment eines Wasserzeichens auf, welches die Harzer\indexO{Harz} Herkunft des Papiers bezeugt. %
Diese Papiersorte ist im Leib\-niz-Nachlass für die Zeitspanne von 1680 bis 1683 bzw. 1687 belegt.
Laut Fachliteratur war das Wasserzeichen ab Oktober 1679 im Umlauf. %
\pend
\pstart
Datierungsrelevant ist ferner die Art, wie \textit{potentia absoluta} und \textit{potentia respectiva} im ersten Teil der Aufzeichnung N.\;\ref{60243}, \refpassage{LH_37_04_087_087r_potabs-1}{LH_37_04_087_087r_potabs-2} voneinander unterschieden werden (siehe für eine Charakterisierung dieses Begriffspaars die ausführlichen Datierungsgründe zu N.\;\ref{58235}, S\;\pageref{58235_dat}).
Dass die bei zusammenstoßenden Körpern anzutreffende "absolute Potenz" -- die skalare Summe ihrer kinetischen Kraft -- als quadratisch ($mv^2$) dargelegt wird, verortet die Aufzeichnung N.\;\ref{60243} chronologisch nach dem Textkomplex \textit{De corporum concursu} (\textit{LBS} VIII,\;3 N.\;58\cite{01401}), den Leibniz eigenhändig auf den Anfang 1678 datiert hat.
Die "relative Potenz" \textendash\ d.\,h. im Kontext der Stoßlehre die \textit{vis percussionis}, die kinetischen Kraft, mit der die zusammenstoßenden Körper aufeinander wirken %
\textendash\ wird indes in dem Fall, dass die Körper mit Geschwindigkeiten zusammenstoßen, die in umgekehrtem Verhältnis zu deren Massen stehen, der Bewegungsgröße (\textit{mv}) gleichgesetzt.
Diese Charakterisierung der \textit{potentia respectiva} ist als Hinweis darauf zu betrachten, dass die Aufzeichnung N.\;\ref{60243} nicht viel später als der Textkomplex \textit{De concursu corporum}\cite{01401} entstand, d.\,h. noch zu einer Zeit, als Leibnizens Kraftlehre wohl noch nicht gefestigt war.
Anders verhält es sich etwa in dem vor Ende November 1689 an Rudolf Christian von Bodenhausen\indexN{\textls{Bodenhausen} (Bodenus, Bodenausen), Rudolf Christian, Freiherr v. 1698} übermittelten Teil von Leibnizens \textit{Dynamica}:
Die "relative Potenz" wird dort stets als quadratisch dargelegt und von der Bewegungsgröße auch dann unterschieden, wenn diese letztere \textendash\ wie in dem besagten Sonderfall \textendash\ beim Stoß erhalten bleiben soll (vgl. \textit{Dynamica}\cite{01354}, pars II, sectio III, prop.\;7--12: %
\textit{CFP}~I, S.\;400--412~\bracks{A}, II.2, S.\;498--512~\bracks{C};\cite{01354} %
siehe über die Entstehung der \textit{Dynamica}\cite{01354} unter besonderer Berücksichtigung des Abschnitts über die Stoßlehre die einleitende Erläuterung zu N.\;\ref{impulsuseccentricus} in diesem Band, \refpassage{impulsecc_intro_Dynam-1}{impulsecc_intro_Dynam-2}).
\pend%
\pstart%
Diese Betrachtung ist Anlass zur Vermutung, dass die Aufzeichnung N.\;\ref{60243} auf den ersten Teil des Belegzeitraums des Wasserzeichens zu datieren ist, d.\,h. auf die Zeitspanne von 1680 bis 1683.
Eine spätere Entstehung (bis 1687) ist jedoch nicht ausgeschlossen.
\pend%
\end{ledgroup}%
\normalsize%
\selectlanguage{latin}%
\frenchspacing%
